\chapter{Respiratory Syncytial Virus (RSV)}

\section*{What Is This Test?}
Respiratory syncytial virus (RSV) testing detects RSV, a single-stranded negative-sense RNA virus of the \textit{Pneumoviridae} family (genus \textit{Orthopneumovirus}). RSV is the most common cause of bronchiolitis and pneumonia in children under 1 year of age and a significant cause of respiratory illness in older adults and immunocompromised patients \cite{hall2009burden}. The virus is transmitted via large respiratory droplets and direct contact with contaminated surfaces, with an incubation period of 2--8 days. Clinical manifestations range from mild upper respiratory symptoms (rhinorrhea, cough) in older children and adults to severe lower respiratory tract disease (bronchiolitis with wheezing, crackles, respiratory distress, and hypoxemia) in infants under 6 months. In older adults (particularly those over 65 with cardiopulmonary disease), RSV causes an illness comparable in severity to influenza, with an estimated 177,000 hospitalizations and 14,000 deaths annually in the United States \cite{falsey2005respiratory}. Transplant recipients and patients with hematologic malignancy are at risk for severe and prolonged RSV disease with high mortality. Two antigenic subgroups (RSV A and RSV B) circulate, often with one predominating in a given season.

\section*{Alternative Names}
\begin{itemize}
  \item RSV Antigen Test
  \item RSV PCR
  \item RSV Rapid Antigen Detection
  \item RSV Culture
  \item RSV DFA
  \item Respiratory Syncytial Virus NAAT
  \item RSV EIA
\end{itemize}
\section*{Principle}
Multiple detection methods are available. Rapid antigen detection tests (RADTs) use lateral-flow immunochromatography to detect RSV fusion (F) protein or nucleoprotein in nasopharyngeal specimens, providing results in 10--15 minutes but with limited sensitivity (50--90\%) compared to molecular methods \cite{chartrand2015diagnostic}. Direct fluorescent antibody (DFA) staining uses fluorescein-labeled monoclonal antibodies to detect RSV antigens in harvested respiratory epithelial cells, with sensitivity of 70--95\%, depending on specimen quality. Reverse-transcription polymerase chain reaction (RT-PCR) detects RSV RNA with sensitivity of 95--100\% and is the gold standard; multiplex panels (BioFire FilmArray, Luminex NxTAG, GenMark ePlex) detect RSV A and B alongside other respiratory pathogens. Viral culture in HEp-2, HeLa, or A549 cells with observation for characteristic syncytia (multinucleated giant cells) is the reference method but requires 3--10 days and has low sensitivity (50--70\%) compared to PCR. Serologic testing by enzyme immunoassay or complement fixation is not used for clinical diagnosis but may be employed in epidemiologic studies.

\section*{History}
RSV was first isolated in 1956 by Robert Chanock from a chimpanzee with coryza (named chimpanzee coryza agent, CCA) and subsequently from infants with bronchiolitis in 1957, when it was renamed respiratory syncytial virus for the characteristic syncytia it produces in cell culture. The rapid diagnosis of RSV via immunofluorescence (FA) was developed in the 1970s. Enzyme immunoassay (EIA)-based antigen detection kits were introduced in the 1980s, followed by rapid immunochromatographic tests in the 1990s. RT-PCR for RSV was developed in the 1990s and became commercially available in the 2000s. FDA approval of palivizumab (Synagis), an RSV-specific monoclonal antibody, in 1998 for prophylaxis in high-risk infants drove demand for rapid and accurate RSV diagnosis. Multiplex respiratory panels emerged in the 2010s. RSV vaccines (Arexvy, Abrysvo) for adults 60 years and older were FDA-approved in 2023, and nirsevimab (Beyfortus), a long-acting monoclonal antibody for infant prophylaxis, was approved in 2023.

\section*{Why This Test Is Ordered}
\begin{itemize}
  \item Diagnosis of RSV infection in infants and young children presenting with bronchiolitis (tachypnea, wheezing, crackles, nasal flaring, retractions) or pneumonia during RSV season (typically October through March) \cite{aap2014bronchiolitis}
  \item Differentiation of RSV from other causes of bronchiolitis (human metapneumovirus, parainfluenza virus, rhinovirus) to guide clinical management and infection control
  \item Evaluation of acute respiratory illness in adults over 65 or with chronic cardiopulmonary disease during winter months, particularly when influenza testing is negative
  \item Surveillance for RSV in immunocompromised patients (hematopoietic stem cell transplant, lung transplant, chemotherapy) who are at high risk for progression to severe lower respiratory tract disease
  \item Infection control decisions in hospitals, including cohorting of RSV-infected patients and implementation of contact and droplet precautions
  \item Qualification of high-risk infants for palivizumab (RSV monoclonal antibody) prophylaxis based on RSV season timing
  \item Epidemiologic surveillance to determine the onset, peak, and end of RSV season, which varies by geographic region
\end{itemize}

\section*{Primary vs Secondary Test}
This is a primary diagnostic test for RSV infection in children under 2 years of age with bronchiolitis during RSV season. In adult and elderly populations, where RSV viral loads are typically lower, RT-PCR is the preferred method. A negative rapid antigen test in a high-risk infant should be confirmed by RT-PCR.

\section*{How This Test Is Performed}
Specimen collection: Nasopharyngeal wash/aspirate is the preferred specimen in infants because it yields a higher cell recovery, improving DFA sensitivity. A small amount of sterile saline is instilled into one nostril and aspirated back into a collection trap. Nasopharyngeal swab (flocked nylon) is an acceptable alternative for older children and adults and for RT-PCR testing. For DFA, the specimen must contain adequate numbers of columnar respiratory epithelial cells. Specimens should be collected within the first 4 days of illness when viral shedding peaks. Transport in viral transport medium at 2--8\textdegree{}C (refrigerated, not frozen) is required; specimens should be processed within 48 hours. For intubated patients, endotracheal aspirate or bronchoalveolar lavage specimens may be tested. Nasal mid-turbinate swabs or combined nasal/OP swabs can be collected from adults.

\section*{What Preparation Is Needed}
No special patient preparation is required. Use of saline or decongestant nasal sprays immediately before collection may reduce specimen cellularity and should be avoided. For infants, suctioning immediately before specimen collection to remove secretions may improve quality. Breastfeeding or bottle feeding immediately before nasopharyngeal aspiration may provoke gagging or vomiting.

\section*{Contraindications and Precautions}
\begin{itemize}
  \item Nasopharyngeal aspiration may cause transient bradycardia in neonates through vagal stimulation; suctioning should be performed with caution and monitoring. Severe coagulopathy or thrombocytopenia (platelets <50,000/$\mu$L) is a relative contraindication. Rapid antigen tests have limited sensitivity (50--90\%), particularly in adults and older children with lower viral loads; negative results should not be used to exclude RSV in high-risk populations. Specimens for culture should not be frozen as this inactivates the virus. False-negative DFA results occur when specimens contain fewer than 20 columnar epithelial cells per low-power field.
\end{itemize}
\section*{Risks}
Nasopharyngeal specimen collection may cause transient discomfort, gagging, coughing, and occasionally mild epistaxis. Nasopharyngeal aspiration in neonates may cause brief oxygen desaturation or bradycardia. All specimen collection should be performed using standard precautions and contact/droplet precautions if RSV is suspected.

\section*{What Happens After the Test}
Rapid antigen test results are available within 10--15 minutes. Rapid molecular (point-of-care) results available in 15--30 minutes. DFA results available in 2--4 hours. RT-PCR results are typically available within 2--6 hours; multiplex respiratory panel results within 1--3 hours. Viral culture results may take 3--10 days. Positive results in hospitalized patients should prompt implementation of contact and droplet precautions. Ribavirin therapy (aerosolized) may be considered for severe RSV disease in immunocompromised patients based on PCR-confirmed diagnosis.

\section*{Understanding Results}

\subsection*{Below Normal}
\begin{itemize}
  \item A negative RSV test (PCR not detected, antigen negative, culture negative) means RSV was not detected in the specimen. This may indicate: Absence of RSV infection; a respiratory illness caused by another virus (rhinovirus, human metapneumovirus, parainfluenza virus, adenovirus, influenza); collection of specimen outside the peak viral shedding period (before 24 hours or beyond 5--7 days from symptom onset); or a false-negative result from rapid antigen testing in a patient with low viral load
  \item In infants with classic bronchiolitis during RSV season, a negative rapid antigen test should be confirmed by RT-PCR, as antigen test sensitivity may be as low as 50\% in some studies
\end{itemize}

\subsection*{Above Normal}
\begin{itemize}
  \item RSV detected by RT-PCR: Positive molecular testing confirms active RSV infection with high accuracy. In infants, this is most commonly associated with bronchiolitis (inflammation of the small airways causing wheezing, hyperinflation, and atelectasis) and/or viral pneumonia. Approximately 2--3\% of infants require hospitalization, and of those hospitalized, 10--20\% require intensive care
  \item RSV detected in adults (typically RSV B positive identified on multiplex respiratory panel): Associated with an influenza-like illness that may progress to pneumonia, particularly in patients over 65 with underlying COPD, congestive heart failure, or immunocompromise. Hospitalized adults with RSV have mortality rates of 6--8\%, comparable to influenza \cite{lee2013high}
  \item RSV detected in immunocompromised patients (HSCT recipients, lung transplant): May indicate progressive lower respiratory tract disease with high mortality (30--60\% in untreated HSCT recipients). Viral loads by quantitative RT-PCR correlate with disease severity and are used to monitor response to therapy
\end{itemize}

\section*{Normal Results}
\textbf{Negative} or \textbf{Not detected} by RT-PCR, antigen testing, or culture. No RSV virus isolated. For DFA: Fewer than 1 RSV-positive cell per well is considered negative.

\section*{Follow-up Tests}
\begin{itemize}
  \item \textbf{Multiplex respiratory pathogen panel (RT-PCR):} When RSV is not detected but clinical presentation suggests viral bronchiolitis or pneumonia; simultaneously tests for RSV A/B, influenza, parainfluenza 1--4, human metapneumovirus, adenovirus, rhinovirus/enterovirus, coronavirus, and \textit{Bordetella pertussis}
  \item \textbf{RSV quantitative RT-PCR:} Monitoring of viral load in immunocompromised patients on therapy (ribavirin); a decline in viral load correlates with treatment response
  \item \textbf{Chest radiograph:} To evaluate for hyperinflation, peribronchial thickening, atelectasis, or consolidation in patients with clinical deterioration despite supportive care
  \item \textbf{Blood gas analysis:} Assessment of hypoxemia and hypercapnia in infants with severe bronchiolitis requiring respiratory support
  \item \textbf{RSV Genotyping (sequencing):} Performed at reference laboratories to identify RSV A versus B and monitor genomic changes for vaccine and immunoprophylaxis development
\end{itemize}

\section*{References}
\bibliographystyle{unsrtnat}
\bibliography{references}

\clearpage
\section*{Regulatory Status and Authorization Date}
This chapter describes a test category, not a specific commercial product. FDA, CDSCO, EU IVDR, and MHRA approval or registration dates therefore cannot be assigned without the assay manufacturer, product name, instrument, and intended use. For laboratory-developed tests, an individual agency approval date may not exist. See the Preface for the interpretation of regulatory status.

\clearpage
